\section{Introduction}

Short-term photovoltaic power forecasting is a productive research area.
Hybrid architectures combining convolutional and recurrent layers, often
with a further decomposition or gradient-boosted stage, appear regularly
in the literature, and reported accuracy improvements are substantial.
A statistical meta-analysis of 180 studies concludes that hybrid models
consistently outperform alternatives and are likely to be the future of
the field [CITE Nguyen Musgens 2021].

We coded 27 papers on hybrid deep learning for photovoltaic power
forecasting, published between 2022 and 2026, on eight dimensions of
evaluation protocol. Coding was conservative: a field was recorded as
not stated unless the paper explicitly stated it, and every judgement is
supported by a verbatim quotation from the source.

Twenty-six of the 27 report no skill score against any reference
forecast. Seventeen compare only against the components of their own
architecture. All 27 report point estimates with no measure of
run-to-run variance. Twenty-two do not state whether their train-test
split preserves temporal order. Seventeen do not state whether night
hours, when power is near zero and trivially predicted, are included in
the reported error.

These are not incidental omissions. Each is a choice that changes the
number a paper reports, and this paper measures by how much.

\subsection{What This Paper Does}

We fix a leakage-controlled evaluation protocol, hold it constant, and
vary individual evaluation choices while keeping the model, the data,
and the training procedure identical. Five architectures are evaluated -
gradient boosting, a recurrent network, a convolutional-recurrent
network, and both recurrent variants with a gradient-boosted residual
correction stage - across three module technologies, three forecast
horizons, two feature regimes and five random seeds, on 900 recorded
runs. The models are instruments for measuring protocol sensitivity
rather than contributions in themselves.

Four results follow.

\textit{The reference forecast changes both the size and the shape of reported
skill.} On one array at a six-hour horizon, the same forecasts score
0.652 against smart persistence and 0.194 against the optimal convex
combination of climatology and persistence recommended by Yang et al.
[CITE Yang 2020] - 70 percent of the apparent skill is attributable to
the reference. More consequentially, skill against persistence rises
monotonically with horizon, the pattern commonly reported as evidence
that a method's advantage grows at longer lead times, while skill
against the proper reference is non-monotonic and peaks at three hours.
The trend is a property of the reference, not the model.

\textit{Night-hour inclusion deflates reported error by a predictable factor.}
Including all 24 hours rather than daylight hours reduces normalised
root-mean-square error by approximately 34 percent at this site, and
the reduction is closed-form: if night errors are near zero, error
scales as the square root of the daylight fraction. Observed and
predicted ratios agree to within 0.3 percent at every array and horizon.
The deflation is a property of latitude and season, not of the
forecasting method. A skill score is nearly invariant to the same
choice.

\textit{How residual training data are constructed can reverse a component's
apparent sign.} Fitting a residual correction stage on validation
residuals and evaluating on validation is in-sample evaluation of that
stage. Doing so turns a component that costs 0.034 in skill into one
that appears to contribute 0.334 - a reversal of sign and an order of
magnitude in size, from a single choice that is not obviously wrong when
described in a methods section.

\textit{Architecture contributes little relative to any of these.} Under a
controlled protocol, gradient boosting and a recurrent network are
statistically indistinguishable at one- and three-hour horizons; the
convolutional front end shows no detectable benefit; and the residual
correction stage reduces skill in all eighteen configurations tested.
The largest difference we measure between any two architectures is 0.024
in skill score. The gap between forecasting with lagged observations and
forecasting with perfect knowledge of future weather is between 0.51 and
0.72 - roughly twenty-five times larger.

All findings hold in direction on a held-out year evaluated once, after
every modelling and evaluation choice was frozen.

\subsection{Contributions}

1. An open, leakage-controlled benchmark harness for short-term
   photovoltaic power forecasting, with explicit alignment assertions,
   documented data exclusions, and 900 machine-readable run records.
   Most tables and figures are generated directly from these records;
   the protocol-sensitivity and significance-testing tables refit models
   from source data, because paired significance testing requires the
   full error series, which the run records do not store. On a
   verification performed on 2026-08-08, all generated outputs were
   deleted and regenerated from a clean state: table files and figure
   previews reproduced byte-identically, and vector figures differed only
   in an embedded creation timestamp.

2. A quantification of four protocol effects - reference forecast,
   night-hour inclusion, residual-stage construction, and weather
   information regime - measured with the model and data held constant,
   and shown to exceed architectural differences by one to two orders of
   magnitude.

3. A component-wise decomposition of the convolutional-recurrent hybrid
   with residual correction, with paired significance testing and seed
   variance, finding no component that justifies its cost on this data.

4. A coded survey of evaluation practice across 27 papers, with verbatim
   supporting evidence, establishing that the choices measured here are
   predominantly unreported in the literature they affect.

\subsection{What This Paper Does Not Claim}

The evidence base is three co-located arrays at a single site, sharing
one weather station and therefore not independent, evaluated over one
held-out year at horizons of one to six hours with deterministic
forecasts. We do not claim that any architecture is generally inferior,
that the protocol effects measured here transfer unchanged to other
sites or climates, or that the surveyed papers are incorrect - only that
their reported numbers cannot be compared with one another, or with
ours, without the protocol information most of them do not provide.

Rigorous evaluation practice exists in solar forecasting and is
documented in detail [CITE Yang 2020], [CITE Yang 2019 ROPES]. One paper
in our sample of 27 follows it [CITE Mayer 2022]. The gap this paper
addresses is not the absence of standards but their non-adoption in a
specific and active subfield.

\subsection{Paper Organisation}

Section II surveys evaluation practice in the hybrid forecasting
literature and situates this work against the solar forecast
verification and machine-learning leakage literatures. Section III
describes the data and preprocessing. Section IV specifies the
evaluation protocol before the models, in that order, because the
protocol is the contribution. Section V gives experimental setup.
Section VI reports results by research question. Sections VII and VIII
give limitations and conclusions.
